\documentclass[11pt]{article}

\usepackage[a4paper,margin=1in]{geometry}
\usepackage{amsmath}
\usepackage{amssymb}
\usepackage{booktabs}
\usepackage{hyperref}
\usepackage{xcolor}
\usepackage{listings}

\lstset{
  basicstyle=\ttfamily\small,
  breaklines=true,
  frame=single,
  columns=fullflexible
}

\title{Context DeLaN: How Object Trajectories Are Predicted}
\author{}
\date{}

\begin{document}
\maketitle

\section{Overview}

The Context DeLaN world model under
\begin{center}
\texttt{scripts/world\_model/Physics/Robot\_and\_Object/Context\_DeLaN}
\end{center}
predicts the manipulated cube's trajectory in a way that is
\emph{structurally different} from its sibling models (MCGDF and
No-Object-Euler) in this directory.  The object side is essentially a
\textbf{black-box acceleration predictor}: one MLP takes the current
state, the recorded torque, and a latent context $z$, and emits the
next-step linear and angular accelerations of the cube.  No Newton's
second law, no Newton--Euler integrator, no contact wrench, no
gripper-twist coupling, and no body-frame inertia tensor appear in the
forward pass.  Everything that happens after the acceleration is
produced is just integration.

This document describes that prediction pipeline end to end, contrasts
it with the explicit-physics versions, and discusses what the
architecture buys and gives up.

\section{State, Input, and Context}

The state is identical to MCGDF and No-Object-Euler:
\begin{equation}
  s_t \;=\;
  \begin{bmatrix}
    q_t \\ \dot q_t \\ p^o_t \\ r^o_t \\ v^o_t \\ \omega^o_t
  \end{bmatrix}
  \in \mathbb{R}^{18 + 13} \;=\; \mathbb{R}^{31},
\end{equation}
the per-step input is the recorded torque
$\tau_t \in \mathbb{R}^{9}$, and an encoder produces a per-episode
latent
\begin{equation}
  z \;=\; \mathrm{ContextEncoder}\bigl(s_{t-K+1:t},\,\tau_{t-K+1:t}\bigr)
  \;\in\; \mathbb{R}^{\mathrm{latent\_dim}}.
  \label{eq:cd-z}
\end{equation}
The encoder is either an LSTM or an MLP (chosen by
\texttt{--context\_encoder lstm|mlp}); both versions live in
\texttt{models.py} alongside the per-step dynamics.

\paragraph{Important property of $z$.}
$z$ is computed once at the start of the rollout and held
\emph{constant} for the full horizon.  There is no online filter on
$z$, no innovation update, no closed-loop refinement -- unlike the
``coupling-driven object dynamics with an online latent context
update'' design described in
\texttt{No\_object\_Euler/robot\_object\_no\_object\_euler\_description.tex}.

\paragraph{Property NOT exploited.}
Although the dataset also records the per-episode physical context
$\xi = [m_o,\, I_b,\, \mu_o]$, the dynamics model never reads it.  It
is loaded for inspection and (optionally) for supervised regression of
$z$, but the forward pass at deployment time depends only on
$(s_t, \tau_t, z)$.

\section{Per-step Object Dynamics: a Single MLP}
\label{sec:cd-object}

The whole object side is the following two lines of code (from
\texttt{ContextRobotObjectDeLaNStep.forward}):
\begin{lstlisting}
coupling_input  = torch.cat([state, torque, z], dim=-1)   # in R^{53}
object_acc      = self.object_acc_net(coupling_input)     # in R^{6}
object_lin_acc  = object_acc[:, :3]   # 3-D linear acceleration
object_ang_acc  = object_acc[:, 3:]   # 3-D angular acceleration
\end{lstlisting}
where \texttt{self.object\_acc\_net} is a 3-layer SiLU MLP built by
\texttt{mlp(state\_dim+torque\_dim+latent\_dim, hidden\_dim=256, 6,
depth=3)}.

In equation form,
\begin{equation}
  \bigl(\dot v^o_t,\,\dot\omega^o_t\bigr)
  \;=\;
  \mathrm{MLP}_{\mathrm{obj}}\!\bigl(s_t,\,\tau_t,\,z\bigr)
  \;\in\;\mathbb{R}^{6}.
  \label{eq:cd-object-acc}
\end{equation}

\paragraph{What is \emph{not} computed here.}
Neither Newton's second law nor the gyroscopic Euler equation is
written down: the network is free to ignore both.  In particular,
\begin{itemize}
  \item the cube mass $m_o$ does not appear anywhere; there is no
    division $F_c / m_o$;
  \item the cube inertia tensor $I_b$ does not appear; there is no
    body-to-world rotation $I_w = R I_b R^\top$ and no
    $\omega^o \times I_w \omega^o$ gyroscopic correction;
  \item gravity is not added analytically; the network has to learn
    from the data that the cube falls under gravity when out of
    contact;
  \item the gripper twist $J(q)\dot q$ is not constructed and does
    not appear in
    Eq.~\eqref{eq:cd-object-acc};
  \item there is no learned contact wrench $F_c$ and therefore no
    structural action--reaction coupling
    $\tau_{\mathrm{contact}} = -J_c^\top F_c$ on the robot side.
\end{itemize}

\paragraph{What auxiliary forces appear in \texttt{aux}.}
The model populates \texttt{aux} with three ``force-like'' tensors
under a unit-mass convention:
\begin{align}
  F^{\mathrm{inertial}}_t &= \dot v^o_t, \\
  F^{\mathrm{gravity}}_t &= (0,\,0,\,-9.81), \\
  F^{\mathrm{external}}_t &= F^{\mathrm{inertial}}_t - F^{\mathrm{gravity}}_t.
\end{align}
These are diagnostic quantities used by the training loss to compare
against the dataset's \texttt{external\_force\_est\_w} field.  They do
\emph{not} enter the forward integrator; they are computed
\emph{from} the predicted acceleration, not the other way around.

\section{Integration: Semi-implicit Euler + Quaternion Exp-Map}

Once the acceleration $(\dot v^o_t, \dot\omega^o_t)$ is in hand, the
object substate is integrated by the same scheme MCGDF and
No-Object-Euler use for the kinematic part:
\begin{align}
  v^o_{t+1}      &= v^o_t      + \dot v^o_t \,\Delta t, \\
  \omega^o_{t+1} &= \omega^o_t + \dot\omega^o_t \,\Delta t, \\
  p^o_{t+1}      &= p^o_t      + v^o_{t+1} \,\Delta t, \\
  r^o_{t+1}      &= r^o_t \otimes \exp\!\bigl(\tfrac{1}{2}\omega^o_{t+1}\,\Delta t\bigr).
\end{align}
Semi-implicit Euler ensures that the new position uses the new
velocity, and the exponential-map quaternion update is the standard
closed-form
\(\exp([0,\mathbf v]) = (\cos\|\mathbf v\|,\,(\sin\|\mathbf v\|/\|\mathbf v\|)\mathbf v)\).
These integrators are analytic (no learnable parameters); the only
learned thing on the object side is the MLP in
Eq.~\eqref{eq:cd-object-acc}.

\section{Trajectory Prediction: a Recursive Rollout}

A predicted trajectory over a horizon $H$ is produced by chaining $H$
calls to the per-step dynamics, with the latent $z$ frozen across all
steps:
\begin{equation}
  s_{t+h+1} \;=\; \mathrm{ContextRobotObjectDeLaNStep}\bigl(s_{t+h},\,\tau_{t+h+1},\,z\bigr),
  \qquad h = 0, 1, \dots, H-1,
  \label{eq:cd-rollout}
\end{equation}
and the predicted object trajectory is the projection
\begin{equation}
  \bigl\{(p^o_{t+h+1},\, r^o_{t+h+1},\, v^o_{t+h+1},\, \omega^o_{t+h+1})\bigr\}_{h=0}^{H-1}
  \;\subset\; s_{t+1:t+H}.
\end{equation}

The corresponding code from
\texttt{ContextRobotObjectDeLaNWorldModel.forward} is:
\begin{lstlisting}
# encode z once, from a K-step history window
z = self.encode_context(history_states, history_torques)

state = history_states[:, -1]
preds = []
for h in range(future_torques.shape[1]):       # H steps
    state, aux = self.dynamics(
        state, future_torques[:, h], z,         # NOTE: z is frozen here
    )
    preds.append(state)
pred_future = torch.stack(preds, dim=1)         # (B, H, 31)
\end{lstlisting}

\paragraph{Where the learned signal flows.}
At step $h$ of the rollout, the object trajectory depends on the
learned model only through Eq.~\eqref{eq:cd-object-acc}: six numbers
per step (3 for $\dot v^o$, 3 for $\dot\omega^o$), repeated $H$ times.
This is structurally similar to MCGDF's 6-D contact-wrench bottleneck
but \emph{interpreted differently}: in MCGDF those six numbers were
$F_c$ and went through Newton--Euler to become accelerations; here
those six numbers \emph{are} the accelerations.

\section{What Is Learned}

\begin{table}[h]
\centering
\small
\begin{tabular}{p{0.30\linewidth} p{0.25\linewidth} p{0.18\linewidth} p{0.16\linewidth}}
\toprule
Sub-network (class) & Input & Output & Notes \\
\midrule
\texttt{ContextEncoder} (LSTM or MLP) & flat history of $(s,\tau)$ over $K$ steps & $z \in \mathbb{R}^{\mathrm{latent\_dim}}$ & Computed once per rollout \\
\texttt{ContextDeLaNStructuredTerms} & $(q, \dot q, z)$ & $H(q,z),\, c(q,\dot q, z),\, g(q,z)$ & Robot-side DeLaN block \\
\texttt{residual\_torque\_net} (MLP) & $(s, \tau, z) \in \mathbb{R}^{53}$ & $\tau_{\mathrm{res}} \in \mathbb{R}^{9}$ & Robot-side residual \\
\texttt{object\_acc\_net} (MLP) & $(s, \tau, z) \in \mathbb{R}^{53}$ & $(\dot v^o,\dot\omega^o) \in \mathbb{R}^{6}$ & \textbf{The whole object side} \\
\bottomrule
\end{tabular}
\caption{Trainable sub-networks of the Context DeLaN model.  Only the
last row produces the object trajectory.}
\label{tab:cd-subnets}
\end{table}

\paragraph{What is \emph{not} learned (object side).}
\begin{itemize}
  \item The semi-implicit Euler integrator for $p^o, v^o, \omega^o$.
  \item The closed-form quaternion exponential-map update for $r^o$.
\end{itemize}
There is no analytic Newton--Euler step on the object side, so the
``physics is in the integrator'' inductive bias that MCGDF and
No-Object-Euler each use in their own way is absent here.

\section{Role of the Latent \texorpdfstring{$z$}{z}}

The latent $z$ is the only mechanism by which Context DeLaN can adapt
to per-episode physical context.  Two consumption sites:
\begin{enumerate}
  \item \emph{Robot side.} $z$ modulates the DeLaN structured terms
    $H(q,z),\,c(q,\dot q,z),\,g(q,z)$ and the residual-torque MLP.
    This is structural FiLM-style conditioning on the robot dynamics.
  \item \emph{Object side.} $z$ is concatenated with $(s,\tau)$ as
    input to \texttt{object\_acc\_net}.  This is plain
    concatenation-conditioning; the network is free to use $z$
    however it likes to shift the predicted acceleration.
\end{enumerate}
Because $z$ is frozen across the rollout, the model relies on the
$K$-step warm-up window to gather enough evidence about the episode
(e.g.\ ``the cube is heavy'' vs ``the cube is light'') to bias the
acceleration head into the right regime.  This is the
\emph{Context-aware Dynamics Model}~/~Pi-CDD lineage's central bet.

\section{Strengths and Limitations}

\paragraph{Strengths.}
\begin{itemize}
  \item \emph{Architectural simplicity.}  The whole object side is
    one MLP plus an analytic integrator.  No quaternion-frame
    bookkeeping for inertia, no $F=ma$ coupling.
  \item \emph{Capacity.}  Because the acceleration head is an
    unconstrained MLP, the model can fit non-rigid-body effects
    (slipping during a pinch, micro-vibrations, soft-finger
    compliance) that the structured MCGDF cannot represent within its
    Newton--Euler integrator.
  \item \emph{Generality across object types.}  With a sufficiently
    informative $z$, the same head can fit heavy and light cubes (and
    in principle anything else in the dataset) without changing the
    architecture.
\end{itemize}

\paragraph{Limitations.}
\begin{itemize}
  \item \emph{No physical inductive bias on the object side.}  The
    network has to relearn $F=ma$, the persistence of gravity, the
    sign of contact forces, and the geometric notion that
    ``displacement to the cube'' should gate motion -- all from data.
    Out-of-distribution behaviour is therefore much harder to
    bound than in MCGDF (where Newton--Euler is enforced) or
    No-Object-Euler (where the carry vs.\ release dichotomy is
    explicit through $\alpha_t$).
  \item \emph{No explicit gripper--object coupling.}  The model has
    to infer from $(q, \dot q)$ that certain gripper configurations
    are ``in contact''.  No $J(q)\dot q$ is provided to the head.
  \item \emph{Frozen $z$ across the rollout.}  The latent cannot be
    refined when an early step's prediction disagrees with what would
    actually happen.  This is exactly what motivates the online
    innovation update proposed in
    \texttt{No\_object\_Euler/robot\_object\_no\_object\_euler\_description.tex}.
  \item \emph{No conservation laws.}  Linear and angular momentum
    conservation between gripper and object is not enforced.
    Equivalently, Newton's third law is not used as a structural
    constraint.
\end{itemize}

\section{Side-by-side Summary}

\begin{table}[h]
\centering
\small
\begin{tabular}{p{0.18\linewidth} p{0.34\linewidth} p{0.20\linewidth} p{0.18\linewidth}}
\toprule
Model & Object-side equation & Learned signal driving the object & Physical prior \\
\midrule
\textbf{MCGDF} & Newton--Euler:
\(m_o\dot v^o = m_o g + F_c^{\mathrm{lin}}\),
\(I_w\dot\omega^o + \omega^o\!\times\! I_w\omega^o = F_c^{\mathrm{ang}}\)
& Contact wrench $F_c \in \mathbb{R}^6$
& Strong: $F=ma$, gyroscopic Euler, world-frame inertia \\
\textbf{No-Object-Euler} & Kinematic blend:
\(v^o_{t+1} = \alpha_t\, v_g + (1-\alpha_t)\,\rho_v\, v^o_t\)
(symmetric for $\omega^o$)
& Coupling factor $\alpha_t \in [0,1]$ (gated by $\|F_c^{\mathrm{lin}}\|$ and $z$)
& Medium: carry vs.\ release dichotomy, no $F=ma$ \\
\textbf{Context DeLaN} & MLP output:
\((\dot v^o, \dot\omega^o) = \mathrm{MLP}(s,\tau,z)\)
& The 6-D acceleration vector itself
& \emph{None.}  No physical structure on the object side. \\
\bottomrule
\end{tabular}
\caption{Comparison of how the three sibling models predict the
object's trajectory.  All three integrate position via semi-implicit
Euler and orientation via the quaternion exponential map.}
\label{tab:cd-vs-others}
\end{table}

\paragraph{Practical implication for evaluation.}
The cross-model prediction-error script
\texttt{compare\_models\_pred\_error\_heavy\_light.py} in this folder
runs all three architectures on the same heavy/light datasets and
reports gripper and cube position error vs.\ rollout horizon.  When
reading those curves, keep in mind that the three models occupy
different points on the physics-prior spectrum: Context DeLaN's
errors reflect the cost of having no structural prior on the object
side, while MCGDF's reflect the cost of being forced to obey
Newton--Euler.

\end{document}
